\documentclass{article}
\usepackage[utf8]{inputenc}
\usepackage[T1]{fontenc}
\usepackage{geometry}
\geometry{
    a4paper,
    left=2.5cm,
    right=2.5cm,
    top=2.5cm,
    bottom=2.5cm
}
\usepackage{graphicx}
\usepackage{float}
\usepackage{hyperref}
\usepackage[english]{babel}
\usepackage{amsmath}
\usepackage{amssymb}
\usepackage{array}
\usepackage{booktabs}
\usepackage{caption}
\usepackage{subcaption}
\usepackage{enumitem}
\usepackage{fancyhdr}
\usepackage{titlesec}
\usepackage{url}
\usepackage{textcomp}
\usepackage{xcolor}
\usepackage{listings}
\usepackage{multirow}
\usepackage{longtable}
\usepackage{tabularx}
\usepackage{times}

% Set up page style
\pagestyle{fancy}
\fancyhf{}
\rfoot{\thepage}
\renewcommand{\headrulewidth}{0pt}

% Title formatting
\titleformat{\chapter}[display]{\normalfont\bfseries}{\chaptertitle\thechapter}{0pt}{\Huge}
\titlespacing*{\chapter}{0pt}{50pt}{40pt}

% Document info
\title{Template Document}
\author{Academic Support Agent}
\date{\today}

\begin{document}
\maketitle

\begin{center}
\includegraphics[scale=0.3]{escudouson.png}
\end{center}

\maketitle

\thispagestyle{empty}

\section{Descripci n General}

\noindent Queremos \textcolor{gray}{insertar el objetivo del proyecto} utilizando \textcolor{gray}{describir los datos principales a utilizar y su origen} para que \textcolor{gray}{describir en qu  sentido va a cambiar el proyecto la forma de resolver un problema o cambios en la organizaci n} y con esto \textcolor{gray}{describir el beneficio principal que se espera con el proyecto.}



\section{Informaci n de contacto}

{\small 
\begin{enumerate}
\item[\faUserGraduate]  \textcolor{gray}{Nombre del estudiante}
\item[\faEnvelope] \textcolor{gray}{correo electr nico}
\end{enumerate}}
\subsubsection*{Director de trabajo terminal}
{\small
\begin{enumerate}
\item[\faUser]  \textcolor{gray}{Nombre del Director de trabajo terminal}
\item[\faEnvelope] \textcolor{gray}{correo electr nico}
\end{enumerate}}
\subsubsection*{Experto asesor \textcolor{gray}{(si aplica)}} 
{\small
\begin{enumerate}
\item[\faUser]  \textcolor{gray}{Nombre del experto asesor}
\item[\faBuilding ]  \textcolor{gray}{Nombre de la instituci n}
\item[\faEnvelope] \textcolor{gray}{correo electr nico}
\end{enumerate}}



\section{Propuesta} 
\noindent \textcolor{gray}{La propuesta debe responder las siguientes preguntas, algunas pueden omitirse si es que no se aplican al proyecto. Proponga subsecciones que respondan las siguientes preguntas, si  corresponde  (NO ESCRIBIR EN FORMATO PREGUNTA - RESPUESTA):  
\begin{itemize}
\item  Qu   problema  se  est   intentando  solucionar?
\item  Qui nes  se  ver n  beneficiados  al  resolver  este problema?
\item Qu   se  est   haciendo  actualmente  para  solucionar  este  problema?
\item  C mo  se  resolver   el  problema  descrito  anteriormente?
\item  Por  qu   la  soluci n que  se propone  es  mejor?
\item  C mo  se  define  el   xito  de  esta  soluci n  y  c mo  se  evaluar   que  la  soluci n  efectivamente resuelve  la  problem tica  planteada  anteriormente?
\item  C mo  se  planea  garantizar  la  precisi n  tecnol gica  /  algor tmica  de  la  soluci n  propuesta  
y c mo  se  mantendr   la  precisi n  continua  de  esta  soluci n?
\item  C mo  se  abordar   la  recopilaci n,  gesti n  y  uso   tico  de  estos  datos,  incluida  la  
seguridad, privacidad,  integridad  y  transparencia  de  los  datos?
\item  Por  qu   tu  equipo  est   en  una  posici n   nica  para  resolver  este  problema  e  implementarla soluci n  propuesta?
\item  Adem s  de  los  miembros  listados  anteriormente,  qui n  m s  es  fundamental  para  el   xito
de  este proyecto?  C mo  se  van  a  construir  estas  estas  colaboraciones?
\end{itemize}}
\newpage


\thispagestyle{empty}

\noindent\textcolor{gray}{Ejemplo nombres de subsecciones:}

\



\subsection*{Planteamiento del problema}

\subsection*{Descripci n de universo de  usuarios}

\subsection*{Soluciones existentes}  

\subsection*{Propuesta y justificaci n de soluci n}

\subsection*{M tricas de evaluaci n}

\subsection*{Posibilidades de escalamiento}

\subsection*{Descripci n de conjunto de datos}

\subsection*{Propuesta de gesti n de datos}

\subsection*{Equipo de trabajo}


\section{Cronograma de actividades}

\noindent \textcolor{gray}{Describir las actividades a realizar el tercer y el cuarto semestre.}

\section{Flujo de ciencia de datos (opcional)}

\noindent \textcolor{gray}{A  grandes  rasgos,  cu l  es  el  proceso  que  se  anticipa  seguir, esquema  para  dise ar,  implementar  y  evaluar  la  soluci n  propuesta.}
%\begin{center}
%    \begin{tikzpicture}[
%        node distance=4ex and 0em,
%        block/.style={rectangle, draw, fill=white!20, 
%    text width=5.5em, text centered, rounded corners, minimum height=4em},
%        line/.style={draw, -latex},
%        ]
%
%        \node [block] (1) {};
%        \node [block, below right= of 1] (2) {};
%        \node [block, below left= of 2] (3) {};
%        \node [block, above left= of 3] (4) {};
%
%        \path [line] (1.east) to[out=0, in=90] (2.60);
%        \path [line] (2.-60) to[out=-90, in=0] (3.east);
%        \path [line] (3.west) to[out=180, in=-90] (4.-120);
%        \path [line] (4.120) to[out=90, in=180] (1.west);
%    \end{tikzpicture}
% \end{center}
\end{document}